% ============================================================
% Chapter 10 — Discussion
% ============================================================
\chapter{Discussion}
\label{ch:discussion}

\roadmap{This chapter interprets the results against the two research
questions, situates them in the literature of \cref{ch:background}, and
examines the threats to validity that condition the interpretation. It is
written to be filled in alongside \cref{ch:results}: where a claim depends
on a study value, it is phrased so that the direction and magnitude come
from the results rather than being asserted here.}

\section{Answering RQ1: can browser sensing detect learner states?}
\label{sec:discussion-rq1}

RQ1 asked whether browser-based multimodal sensing can detect learner
affective and cognitive states with useful fidelity, validated against
self-report. The cross-modal analysis (\cref{sec:results-validation}) and
the self-report agreement (\cref{fig:validation-confusion}) bear on this
directly. \todo{Interpret the observed $r$ and $\kappa$: if agreement is
moderate, argue that facial and gaze channels converge on engagement at
the region level, consistent with the modest multimodal gains reported by
\citet{dmello2015review}; if agreement is weak, attribute it to the shared
webcam bottleneck and region-level gaze, and to face-loss rates
(\cref{tab:results-yield}) of the kind \citet{bosch2015} documented in the
wild.} Three interpretive commitments hold regardless of the numbers.

First, the validation is \emph{convergent}, not criterion
\citep{campbell1959}: agreement between the facial and gaze channels
raises confidence that both track a common engagement construct, but
neither is ground truth, and correlated error (poor lighting degrading
both) inflates apparent agreement. The self-report probe partially breaks
this dependence---it is a different method---but is sparse and, for the
legacy sessions, absent.

Second, the state space is not uniformly detectable. The affect-dynamics
account \citep{dmello2012dynamics} and the wild-detection literature
\citep{bosch2016} both predict that engagement and boredom are more
separable from faces than confusion and frustration, which share
action-unit signatures (brow lowering, lid tightening). \todo{Report which
states showed the strongest and weakest agreement and relate them to this
prediction.}

Third, the linguistic channel is complementary. Where facial signals are
ambiguous, the EF text-mining detections (\cref{sec:results-ef}) provide
an independent read on confusion, frustration, and metacognitive
monitoring \citep{winne1998}, and the report treats convergence between
the facial and linguistic channels as additional---if still
non-criterion---evidence.

\section{Answering RQ2: does the platform support learning?}
\label{sec:discussion-rq2}

RQ2 asked whether the LLM-grounded platform with learner-initiated
strategy interventions supports learning. The pre/post gain
(\cref{sec:results-gain}) and the intervention-engagement data
(\cref{sec:results-interventions}) bear on this. \todo{Interpret the
normalized gain and its effect size: relate the magnitude to the
laboratory effects for retrieval and spaced practice \citep{roediger2006,
cepeda2006, dunlosky2013}, noting that learner-initiated, in-the-wild use
should be expected to yield smaller and more variable effects than
experimenter-controlled administration.}

The learner-initiated design reframes the pedagogical question as a
self-regulation one. Whether learners \emph{elect} to use effective
strategies, and whether elected use associates with gain, speaks to
autonomy-supportive engagement \citep{ryan2000} more than to the
strategies' intrinsic efficacy, which the cited literature already
establishes. \todo{Report suggestion-uptake and completion-versus-dismissal
patterns and interpret them as evidence about learner self-regulation and
the acceptability of the interventions.}

Because interventions are never affect-triggered
(\cref{sec:interventions-autotrigger}), RQ2's result is uncontaminated by
detection error---a deliberate methodological gain that also caps the
achievable effect, since the platform cannot yet direct support to the
moments of confusion where it would help most. That capping is the bridge
to \cref{ch:future}.

\section{The measurement instrument as contribution}
\label{sec:discussion-instrument}

Independent of the two questions' quantitative answers, the platform
demonstrates that a complete LLM tutoring environment can be instrumented,
on consumer hardware, to a depth that supports replay, retrospective
coding, and multi-channel validation. The AOI attention layer contributes
a task-referenced attention measure that raw gaze heatmaps lack
(\cref{sec:sensing-aoi}), and the decoupling of sensing from pedagogy
gives a clean observational base for detection research. This
instrument---not any single detection accuracy---is the project's most
durable output.

\section{Threats to validity}
\label{sec:discussion-threats}

\paragraph{Construct validity.} Engagement is measured as a facial and
behavioral estimate, not as the latent state itself; the report's
consistent ``estimate/proxy/reference criterion'' language is a guard
against reifying it. The pupil channel's construct validity is low by
construction (\cref{sec:sensing-pupil}) and is excluded from primary
claims. The two low-feasibility EF constructs are flagged in the data and
in reporting.

\paragraph{Internal validity.} The absence of a phase marker
(\cref{sec:methodology-stats}) means any group comparison rests on an
externally defined grouping; the pre/post operationalization is defensible
but is not a randomized manipulation, so gain cannot be attributed to any
single platform feature. The data-quality caveats
(\cref{sec:sensing-caveats})---frame truncation, silent batch drops,
docked-scroll gaps---introduce missingness that is plausibly not at random
(longer or more interactive sessions are differentially affected), and the
analysis must report how affected sessions were handled.

\paragraph{External validity.} A single-cohort study on specific course
material with \todo{N} participants limits generalization; webcam quality,
lighting, and demographics differ from any target population, and
WebGazer's accuracy is device-dependent.

\paragraph{Statistical-conclusion validity.} Small $n$ limits power for
the group comparisons and for per-state agreement; the report accompanies
every test with an effect size and confidence interval
\citep{cohen1988} rather than relying on significance alone, and treats
non-significant results as inconclusive rather than null.

\medskip
\noindent These threats motivate the explicit limitations catalogue of the
next chapter, which states each constraint as a bounded scientific claim
rather than an apology.
